\section{Introduction}\label{sec:intro}

\IEEEPARstart{T}{he} pervasive deployment of Internet-of-Things (IoT) devices across smart cities, industrial control, healthcare, and autonomous transport has dramatically expanded the network attack surface. Distributed denial-of-service (DDoS) campaigns, botnets, reconnaissance scans, and other intrusions traverse the same edge gateways that carry legitimate IoT traffic, motivating the deployment of network intrusion detection systems (IDS) at the network edge \cite{nguyen2021fl-iot-survey}.

\textbf{Federated learning for IoT IDS.} Centralised IDS pipelines, in which raw flow records are exfiltrated to a central server for model training, raise persistent privacy concerns and conflict with regulations such as GDPR and HIPAA. Federated Learning (FL) \cite{mcmahan2017fedavg} offers a privacy-preserving alternative: each IoT gateway trains a local model on its own traffic and only the model updates are aggregated. A growing literature applies FL to IoT IDS, with notable contributions including DeepFed \cite{li2020deepfed}, which combined CNNs with gated recurrent units, and a wave of subsequent work on heterogeneity-robust aggregation, communication compression \cite{sattler2019sparse,reisizadeh2020fedpaq,aji2017topk}, and security-aware client selection.

\textbf{The communication--robustness tension.} Two pressures shape FL-IDS design at the edge. First, IoT links are bandwidth-constrained: a typical LoRaWAN device has a duty-cycled uplink budget on the order of kilobytes per hour. Second, client data is highly non-identical. The geographically distributed nature of IoT deployments means that one gateway may see almost only DDoS traffic while another sees primarily reconnaissance scans, a scenario commonly modelled by a Dirichlet partition with a small concentration parameter $\alpha$ \cite{hsu2019measuring}. Under such partitioning, vanilla FedAvg-MLP architectures are known to suffer slow convergence and, more pathologically, occasional catastrophic divergence in which a single unfortunate client--initialisation pair drags the global model far from the consensus optimum.

\textbf{Kolmogorov--Arnold Networks (KANs).} The recent introduction of KANs \cite{liu2024kan} replaces fixed neuron-wise activations with learnable edge-wise B-spline activations. Unlike a dense MLP weight matrix, where each scalar parameter modulates a global linear interaction, each spline coefficient in a KAN layer has \emph{local} support on the grid by construction of the B-spline basis. This locality has been argued to provide better interpretability and parameter efficiency in physics-discovery and scientific-computing settings.

\textbf{KAN-FL: a recently busy research line.} Within a year of \cite{liu2024kan}, several works have integrated KANs into federated learning. F-KANs \cite{zeydan2024fkans} ported the original KAN to a vanilla FedAvg framework on CIFAR-style benchmarks. \cite{sasse2024evaluating} performed the first systematic non-IID evaluation of FedKAN, reporting modest gains over MLPs at matched parameter budgets on tabular data. \cite{ma2025enhancing} extended the comparison across multiple aggregation strategies (FedProx, SCAFFOLD), and \cite{zeydan2025fedkan} applied a Fed-KAN to traffic prediction. \cite{lee2025unified} contributed a medical-imaging benchmark and \cite{zeleke2024federated} an ECG-signal study. Most recently, \cite{yu2025cgfkan} proposed CG-FKAN, which sparsifies the KAN grid coefficients to attack communication efficiency directly.

\textbf{Positioning of this paper.} The methodological observation ``KAN-FL converges in fewer rounds and transmits fewer bytes than a parameter-heavy MLP'' has therefore been published several times outside the IDS setting. We do \emph{not} claim that observation as new. What is new in this paper is (i) a systematic empirical study of FedKAN \emph{specifically} on the standardised NetFlow-v2 family of IDS datasets~\cite{sarhan2022standard}, (ii) a parameter-matched comparison protocol that resolves the unfair-baseline ambiguity present in much prior FedKAN work, and (iii) a refined narrative anchored in a convergence analysis: FedKAN's advantage is not in mean accuracy but in worst-case robustness under extreme client heterogeneity, and this advantage cannot be replicated by simply scaling up an MLP baseline.

\textbf{Headline empirical finding.} On NF-BoT-IoT-v2 binary detection under Dirichlet $\alpha = 0.1$ (extreme non-IID) with $K=10$ clients and $T=50$ rounds, averaged over $n=10$ seeds: the parameter-matched FedAvg-MLP-PM (3,362 parameters) collapses to as low as 56.5\% accuracy on its worst seed, while FedKAN-8 (3,280 parameters) maintains $\geq 80.7\%$ on every seed, with a $2.6\times$ smaller cross-seed standard deviation. The seed-paired bootstrap 95\% confidence interval for the mean accuracy difference is $[+0.50, +14.57]$ percentage points. Doubling MLP capacity to 3.4k parameters \emph{worsens} the worst-seed accuracy from 60.4\% (small MLP at 1.3k params) to 56.5\%, ruling out a capacity-driven explanation.

\textbf{Theoretical anchoring.} We provide a convergence analysis (Theorem~\ref{thm:conv}) showing that the FedKAN error decomposes into a standard FedAvg term proportional to the data-heterogeneity factor $\Gamma_\alpha$ \emph{plus} an additive spline-approximation bias of order $\mathcal{O}(G^{-(k+1)})$. Crucially, the spline-bias term is independent of $\Gamma_\alpha$, behaving as an effective regulariser whose relative impact \emph{decreases} as heterogeneity grows. Empirically, this manifests as the smaller cross-seed variance we report.

\textbf{Contributions.}
\begin{itemize}
\item We present the first systematic, parameter-matched empirical comparison of FedKAN against FedAvg-MLP on the standardised NetFlow-v2 IDS datasets \cite{sarhan2022standard}, under both binary and 5-class detection and three data-heterogeneity regimes (IID, Dirichlet $\alpha\in\{1.0, 0.1\}$). Our protocol explicitly addresses the parameter-matching gap identified in prior FedKAN work.
\item We identify a robustness/tail-risk advantage of FedKAN that is architectural rather than capacity-driven: parameter-matched MLP suffers worst-seed accuracy drops of $\sim 24$ percentage points relative to FedKAN-8 under extreme non-IID, and doubling MLP capacity does not close this gap. The advantage is supported by a seed-paired bootstrap 95\% CI excluding zero across $n=10$ seeds.
\item We anchor the empirical finding in a convergence analysis introducing a spline-bias term $\mathcal{O}(G^{-(k+1)})$ that is independent of the heterogeneity factor and acts as an effective regulariser (Theorem~\ref{thm:conv} and Corollary~\ref{cor:robust}).
\item We honestly walk back the ``$\sim 50\%$ communication reduction'' framing of prior FedKAN-IDS work, showing via Proposition~\ref{prop:comm} that under parameter parity FedKAN incurs roughly $14\%$ \emph{more} bytes per round than its MLP counterpart, and reframe the case for FedKAN around tail-risk reduction, the more deployment-relevant metric for IDS.
\item All code, dataset preparation scripts, configurations, and $289$ reproducible runs are released at \url{https://github.com/haodpsut/FedKAN-IDS-v2}.
\end{itemize}

\textbf{Organisation.} Section~\ref{sec:prelim} introduces the KAN layer formalism and the federated IDS optimisation problem. Section~\ref{sec:method} details the FedKAN-IDS framework, including server and client algorithms. Section~\ref{sec:theory} presents the four formal results that anchor our empirical narrative. Section~\ref{sec:experiments} reports experimental protocol and results across architectures, partitions, and modes. Section~\ref{sec:conclusion} discusses threats to validity and outlines next steps.
